\clearpage
\section{Prime-Shell Signals and the Frey Nonlinear Transform}

\subsection{Prime support as bandwidth and valuation as persistence}
For
\[
 n=\prod_p p^{e_p},
\]
define the prime-shell signal
\[
 \Psi_n(p,r)
 =
 \sqrt{\log p}\,\mathbf 1_{0\le r<e_p}.
\]
Its total energy is
\[
 \norm{\Psi_n}_2^2
 =
 \sum_p e_p\log p
 =
 \boxed{\log n}.
\]
Let $P_{\mathrm{act}}$ retain only the first shell in each active prime channel. Then
\[
 \norm{P_{\mathrm{act}}\Psi_n}_2^2
 =
 \sum_{p\mid n}\log p
 =
 \boxed{\log\rad(n)}.
\]
Thus the radical is activated bandwidth, while the difference
\[
 \log n-\log\rad(n)
\]
is repeated persistence depth.

For a primitive triple, the direct-sum signal
\[
 \Psi_{a,b,c}
 =
 \Psi_a\oplus\Psi_b\oplus\Psi_c
\]
has disjoint prime support in the three sectors. The standard quality becomes the exact energy ratio
\[
 \boxed{
 q(a,b,c)
 =
 \frac{\log c}{\log\rad(abc)}
 =
 \frac{\norm{P_c\Psi_{a,b,c}}_2^2}
 {\norm{P_{\mathrm{act}}\Psi_{a,b,c}}_2^2}.
 }
\]
High quality means large output persistence relative to activated global bandwidth.

\subsection{Unitary transforms cannot create a radical saving}
Apply a unitary Fourier, wavelet, or polyphase transform in the shell variable. Parseval preserves total shell energy. Therefore the ratio between total depth energy and activated support energy is unchanged by a unitary shell analysis followed by exact synthesis.

\begin{proposition}[Unitary no-gain principle]
Let $U$ be unitary on the shell coordinate of every active prime channel. Then
\[
 \norm{U\Psi_n}_2=\norm{\Psi_n}_2.
\]
Transport the activation projector to $UP_{\mathrm{act}}U^{-1}$. Its projected energy then also agrees exactly with the original activation energy. The same holds after exact synthesis. An arbitrary unitary transform need not preserve the energy under a fixed, untransported first-shell projection.
\end{proposition}

Signal theory is still useful because it reveals which nonunitary or nonlinear operations would be required. It also separates local persistence from global support before introducing a nonlinear arithmetic carrier.

\subsection{The three-spike \texorpdfstring{$p$-adic}{p-adic} coherence signal}
Represent the additive triple by
\[
 \nu_{a,b,c}=\delta_0+\delta_a+\delta_c.
\]
The pairwise differences are $a$, $b$, and $c$. At a prime dividing exactly one of these integers, exactly one pair of spikes is $p$-adically close. The depth of that edge is the corresponding valuation.

The triple therefore determines a prime-indexed coherence triangle. The radical counts the active triangles. The valuations count how long a pair remains coupled through the $p$-adic refinement tree. This picture anticipates the later cross-prime logarithmic network.

\subsection{The Frey curve}
Attach the marked Frey curve
\[
 \boxed{
 E_{a,b,c}:y^2=x(x-a)(x+b).
 }
\]
The roots $0$, $a$, and $-b$ retain the three differences $a$, $b$, and $c$. The basic invariants are
\[
 \Delta_E=16a^2b^2c^2
\]
and
\[
 c_4(E)=16(a^2+ab+b^2).
\]
At every odd prime $p\mid abc$, the curve has multiplicative reduction. The conductor records that the channel is active, while the minimal discriminant records the full depth.

In signal language, the Frey curve is a nonlinear encoder that converts:
\[
 \begin{aligned}
 \text{prime support}&\longrightarrow\text{conductor},\\
 \text{valuation depth}&\longrightarrow\text{minimal discriminant depth},\\
 \text{additive relation}&\longrightarrow\text{elliptic geometry}.
 \end{aligned}
\]
This is precisely the kind of nonlinear transform that a unitary shell system could not provide.

\subsection{The automorphic eigenchannel}
Modularity associates a weight-two newform $f_E$ to the rational isogeny class of $E$. Its Hecke coefficients
\[
 a_\ell(E)
\]
form a global frequency signal. The conductor is the automorphic bandwidth. The newform, however, remembers the rational isogeny class rather than the marked integral representative.

This information loss is essential. Different integral Frey models can lie in the same rational isogeny or twist channel while having different minimal discriminant depth, N\'eron differential, and modular synthesis cost. A proof cannot use the rational eigen-signal alone.

\subsection{Dyadic Landen descent}
Let $u>v>0$ be coprime odd integers. Consider
\[
 T_0=(v^2,u^2-v^2,u^2)
\]
and
\[
 T_1=
 \left(
 uv,
 \left(\frac{u-v}{2}\right)^2,
 \left(\frac{u+v}{2}\right)^2
 \right).
\]
Both are primitive triples. Their radicals are equal:
\[
 \rad(T_0)=\rad(T_1)=\rad\bigl(uv(u-v)(u+v)\bigr).
\]
The second output is strictly smaller:
\[
 \left(\frac{u+v}{2}\right)^2<u^2.
\]
Hence the second quality is smaller.

The Legendre parameter of the parent is
\[
 \lambda=t^2,
 \qquad
 t=v/u.
\]
The degree-two Landen transformation gives
\[
 \boxed{
 \lambda
 \longmapsto
 \frac{4\sqrt\lambda}{(1+\sqrt\lambda)^2}.
 }
\]
This is exactly the child ratio $a_1/c_1$. The transformation preserves the automorphic channel up to the corresponding two-isogeny while lowering the integral height. It proves that the integral representative inside an isogeny class is relevant to $abc$.

\subsection{Exact automorphic analysis--synthesis identity}
Let
\[
 \phi_E:X_0(N_E)\longrightarrow E
\]
be a nonconstant modular parametrization, and let
\[
 \phi_E^*\omega_E
 =
 2\pi i\,c_\phi f_E(z)\,dz.
\]
Here $c_\phi$ is the differential multiplier of this chosen parametrization and differential. The marked Frey curve need not be the optimal representative of its isogeny class. Define the differential area
\[
 A(E)
 =
 \frac{i}{2}
 \int_{E(\C)}
 \omega_E\wedge\overline{\omega_E}.
\]
Pulling the area form back gives
\[
 \boxed{
 \deg(\phi_E)A(E)
 =
 4\pi^2 c_\phi^2
 \norm{f_E}_{\mathrm{Pet}}^2.
 }
\]
For the Frey model, scaling $x=cX$, $y=c^{3/2}Y$ gives an area factor proportional to $1/c$. Thus the height $c$ can be written as a shape-corrected synthesis gain involving modular degree, Manin normalization, and Petersson energy.

The identity does not prove a conductor-scale upper bound. It identifies the exact quantities that an integral automorphic frame theorem would have to control.

